\section{Further Misunderstandings of the Pagan Worldview}

Evola continues his critique by pointing out that the goal should be to transcend Christianity by integrating it into something more comprehensive that incorporates the essential elements of the genuine antique paganism. This is actually \emph{analogous} to the Catholic position as described by Fr. \textbf{Heinrich Denifle}:

\begin{quotex}
The Catholic Church is the religion. As the supernatural does not destroy the natural, but rather presupposes it, so the Catholic Church does not annul the natural religion but embraces it as a necessary supposition and foundation.
\end{quotex}


Thus, Christianity as understood by its ablest apologists is not averse to embracing the authentically Traditional elements of the earliest religions including the Vedanta and Western paganism.

Once all that has been established, the possibility we suggested of ``transcending'' certain aspects of Christianity would then become real. According to its strict Latin etymology, to transcend means to surpass while ascending. In principle, it is not a matter -- it is good to repeat -- of rejecting Christianity or of showing the same incomprehension towards it as Christianity itself has shown, and largely continues to show, towards paganism. It would instead perhaps be a matter of integrating it into something much larger, omitting some aspects, for which it is little in accord with the spirit typical of some of the recent transalpine innovations in order to accentuate instead other more essential aspects by which such faiths cannot contradict the general conceptions of Aryan and Nordic, pre-Christian and non-Christian spirituality.

\begin{quotationx}
Unfortunately, the neo-pagans fabricate a sham spirituality that rejects the Traditional elements of Christendom while recovering little from genuine antique paganism. He points out the anti-Catholicism of the neo-pagans links them to the elements of the anti-Tradition. As for the historical references: the 1870s was the time of the unification of Italy under a secular Republic on the French model. The Fascists, on the other hand, declared Catholicism the official religion of Italy and restored religious education to the public schools. Evola seems to have moderated his view (in 1941) from the views expressed in Pagan Imperialism (1927), some of which he regrets. Instead of the return to the antique religion of the pagans, what he witnessed was the rise of neo-paganism. More than once, he conceded that those neo-pagans would be better off as Catholics.

Evola is in an odd position: while devoid of any faith and lacking any sympathy with Christian belief, he nevertheless defends Christendom, even referring to it as the time when well-born men held sound beliefs. In this paragraph, for example, he opposes those the Church regarded as ``rebels and heretics''. I would recommend Revolution and Counter-Revolution\footnote{\url{http://www.pliniocorreadeoliveira.info/UK\_RCR.pdf}} for a reading of history that, apart from its theological orientation, has many points in common with Evola. Specifically, they both see the Renaissance, the Reformation, the Enlightenment, Revolution, democracy, and so on, not as the Progress that most believe them to be, but rather as causes of the end of a phase of Tradition.
\end{quotationx}

Unfortunately, as we pointed out, the path taken by the extremist racist neo-paganism is totally different. Falling into a snare almost deliberately set, these neo-pagans, end up by advocating and defending ideas more or less reducing them to that sham, naturalistic, particularistic paganism, lacking light and transcendence, and even pervaded by a misunderstood, pantheistic mysticism, which was created polemically from the Christian misunderstanding of the pre-Christian world, and which is based, at most, only on isolated forms of degeneration of the regressions of such a world. And as if this were not enough, they often resort to an anti-Catholic polemic which, \emph{mutatis mutandi}, whatever its political justification, in fact drags out certain arguments and clichés of a purely modern, rationalist, protestant, and enlightenment type that were the weapons of liberalism, democracy, and freemasonry. This was already the case, to a degree, with H. S. Chamberlain. It appears again in a certain Italian racialist movements, which are inspired by Gentilian philosophy, that is, the philosophy of personalism through which Fascism would be the continuation of the anti-Catholicism of the 1870s, Roman re-enactments would be a silly rhetoric, and the Italian tradition would tend to coincide more or less with the opinions of a series of rebels and heretics, Giordano Bruno until now.

Generally, what we just indicated is recognized when the new paganism gives itself up to the glorification of immanence, ``Life'', ``nature'', seeking to create new, superstitious mysticism, which is in the sharpest contrast to every higher Olympian and heroic ideal of the great Aryan civilizations of pre-Christian antiquity. What are we to think about assertions like this one? ``Faith in a super-sensible world beyond this world of sense is something schizophrenic: only schizophrenics see double''\footnote{E. Bergmann}. Or about another, according to which every distinction between spirit and body is only a product of the anti-Aryan degeneration injected by oriental races.

\begin{quotex}
Evola defends the immortality of the soul. He mocks the idea of immortality understood as a genetic continuation, an idea that today is expressed in concepts like the selfish gene or human biodiversity.
\end{quotex}

Denying this distinction, such racialists, as a consequence, come to deny immortality itself: if the soul is inconceivable separately from the body, one cannot think of a survival in the afterlife, but only of immortality meant as self-continuation through one's descendants. An immortality, which a massacre, earthquake, or epidemic naturally enough would be sufficient to destroy.

\begin{quotex}
Whatever his admiration for Nietzsche may have been, Evola shows he is a rather inconsistent Nietzschean. Evola consistently criticizes Nietzsche's incomprehension of the ascetic life.
\end{quotex}

Regarding the anti-ascetic prejudice, we already said: neo-paganism redoubles, in this regard, the incomprehension already demonstrated by Nietzsche: The Aryan would not have known ascesis in normal life: his true mysticism would have been for the afterlife: it would never have thought of a supernatural fulfillment of the personality.

\begin{quotex}
Evola points out the importance of rite and ritual in the \textbf{Ancient City}. Evola points to the example of infallibility to demonstrate the untraditional nature of the neo-pagans. Many claim to be anti-egalitarian in theory, but in practice they chafe at the notion that one man's thoughts may be superior to another's.
\end{quotex}

Superstitions, residues of the ``dark ages'' and ``Etruscan magic'', lies and instruments for the tactic of temporal control by the clergy in the commerce of indulgences would be, for the others, everything that is sacrament and rite and supernatural power. Thus, they show they do not know that all life in those ancient Aryan and specifically of the Roman pagan civilizations, always had a ritual character, the rite accompanying every form of individual and collective life, not in the sense of an empty ceremony, but as the instrument of a real connection between the human and super-sensible forces. On the other hand, Chamberlain already went to attribute to the Aryan spirit the achievements typical of so-called free enquiry and modern profane science. Naturally, when one believes that Lutheranism represented the awakening of the spirit of the Nordic race instead of representing, as it did represent, the incentive for its further spiritual involution and for its eventual semitization. Elsewhere, in the German edition of Revolt against the modern world, we justified this view to an incomprehension that can only be added to incomprehension. So there is something ingenuous, as Guenon correctly pointed out, in the scandal that was demonstrated among the protestants in the face of the claim of infallibility in the order of transcendent knowledge (in the West it is said: in matters of ``faith''); for the ancient Aryan civilization infallibility was instead recognized as belonging not to a single man, as in Catholicism, but to every member belonging legitimately to the brahman, to the ``solar caste'' of spiritual leaders.

\begin{quotex}
Evola mentions the possibility of a return to antique traditions. However, this poses an insurmountable obstacle even to those pagans today who understand this critique and to whom it may not apply. There is no way today to reconstruct those rites and rituals. Furthermore, their source is of divine origin and tied to their ancestors, blood, and soil. That cannot be done \emph{di nuovo} by fiat without some divine sanction; in particular, one cannot be a pagan by oneself apart from an hierarchically organized community led by a high priest. Those ``sacred and spiritual origins'' are lost in time.
\end{quotex}

In the face of such confusions, we can always put very clearly the alternative: either return to the traditions and their sacred and spiritual origins, or continue to gamble with various combinations and inclinations of modern and secular thought. Another example: what is that ``nature'' that in certain racialist circles is so exalted? It would little suffice to notice that it is not in the least the nature that ancient man saw, but a rationalist construction of the period of French Encyclopedism. The encyclopedists created, with precise subversive and revolutionary intentions, the myth of a good, wise and provident nature in contrast with the corruption of every ``culture''. Thus, we see that the optimistic naturalist myth of a Rousseau and the encyclopedists went hand in hand with ``natural right'', universalism, humanitarianism, egalitarianism, and the denial of every positive form of State and hierarchy.

\begin{quotex}
Evola criticizes the very notion of a return to ``nature''. This applies \emph{a fortiori} to contemporary and subtle variations based on zoological racism. The intent of Sintesi is to develop the spiritual notion of race, not to defend any biological understanding of it.
\end{quotex}

Even in regard to the natural sciences one could say the same. Every honest scientist knows that in his researches, whose goal exclusively consists of the formulation of abstract uniformities and mathematical relationships, there is no place for ``nature''. As for biological research, the very science of genetics, and so on, we already emphasized the errors and the one-sidedness in which they fall, in believing definite laws that are of value only as a partial and subordinate aspect of reality. There is not a trace of what people today consider to be ``nature'' in the meaning that nature had for the man of the origins, for the traditional and solar man characterized essentially by his Olympic and regal distance. Since Italian racialism has not yet moved into those area, it is therefore well to pay attention, and as we said, to value someone else's experience.


\flrightit{Posted on 2011-07-27 by Aeneas}

